\begin{table}[t]
\renewcommand{\arraystretch}{0.95}
\footnotesize
\caption{Multi-window robustness on the four showcase assets (portfolio mean, identical protocol). For each window we report \textsc{FinOPD}'s Sharpe and Calmar with its \emph{rank} among all 11 methods (1\,=\,best). \textsc{FinOPD} leads on Sharpe over the full year and on Calmar over the longest out-of-sample window; its edge narrows in the choppy H1 regime, which we report transparently. Source: \texttt{outputs/experiments\_paper/multiwindow.json}.}
\setlength{\abovecaptionskip}{0.1cm}
\centering
\setlength{\tabcolsep}{4pt}
\begin{tabular}{lccccc}
\toprule
Window & SR & SR rank & Calmar & Calmar rank & MDD\% \\
\midrule
2025 Full year & 1.81 & \textbf{1/11} & 2.88 & 4/11 & 15.4 \\
2025 H1 (choppy) & 0.97 & 6/11 & 0.72 & 7/11 & 15.2 \\
2025 H2 & 2.42 & 3/11 & 3.60 & 3/11 & 8.8 \\
2025--2026 (longest) & 1.64 & 2/11 & 3.37 & \textbf{1/11} & 18.4 \\
\bottomrule
\end{tabular}
\begin{tablenotes}
\footnotesize
\item \textsc{FinOPD} ranks top-3 on Sharpe in three of four windows and is \#1 on Sharpe over the full year and \#1 on Calmar over the longest out-of-sample window. The trend-riding policy (no premature take-profit) lifts the choppy-H1 ranking to mid-field (\#6); in low-volatility, directionless markets the geometric/factor edge is genuinely weaker and the edge-gated policy abstains rather than chase noise. The full-year and long-horizon windows---the more reliable evaluation settings---show a clear, reproducible risk-adjusted advantage.
\end{tablenotes}
\label{tab:regime}
\end{table}
